\subsection{US Replication Design}
\label{sec:methods-us}

We first combined section GPAs into one mean for each university, broad area and year, weighting sections by total enrolment, as in Jeong et al.\ (``Construction''). This gave us 1,845 cells from 29 institutions. US departments do not match the 29 Spanish fields directly, so we grouped both countries into ten common areas: Education; Arts, media and architecture; Humanities and languages; Social sciences and law; Business; Science, mathematics and computing; Engineering; Agriculture; Health and veterinary; and Sports, tourism and services. We used the authors' department workbook and ordered name-matching rules, leaving out ambiguous departments. The mappings can be reviewed in \texttt{data/mappings/}. For Spain, we averaged the available field means equally within each broad-area cell, giving 5,691 cells from 89 universities.

We then fitted separate ordinary least-squares models with area, university and year fixed effects, giving each cell equal weight. Sum coding makes the ten area effects add to zero. We compared their signs and ordering using Pearson and Spearman correlations, with a Fisher-$z$ 95\% interval for Pearson correlation. Figure~\ref{fig:us} uses coefficient intervals clustered by university. We also checked equal section weights across all 33 institutions, A+ capped at 4.0, common years (2015--2023), and leaving out each US university. We compare relative patterns because the grade scales differ (0--4.3 versus 5--10), as do section and graduate averages. Spain also includes private universities.
